% Part: lambda-calculus
% Chapter: lambda-definability
% Section: partial-recursive-functions

\documentclass[../../../include/open-logic-section]{subfiles}

\begin{document}

\olfileid{lam}{ldf}{par}
\olsection{جزوی بازگشتی فنکشن \usetoken{S}{lambda definable} نیں}

جزوی بازگشتی فنکشن اوہ نیں جیہڑے مُڈھلے فنکشوناں توں
ترکیب، بنیادی بازگشت تے بغیر حد گھٹ توں گھٹ حل لبھن نال
حاصل ہوندے نیں۔ ایہ عام بازگشتی فنکشوناں توں ایس گل وچ
وکھرے نیں کہ بغیر حد دی تلاش وچ ورتے فنکشوناں دا باقاعدہ
ہونا ضروری نہیں۔ باقاعدگی دی شرط نہ ہووے، تاں گھٹ توں
گھٹ حل لبھن نال متعین فنکشن کدے غیر معین وی ہو سکدا اے۔

پہلی نظر وچ لگ سکدا اے کہ جیہڑے طریقیاں نال (ہر انپٹ
اُتے معین) سارے عام بازگشتی فنکشن !!{lambda definable}
دکھائے سن، اوہناں نال سارے جزوی بازگشتی فنکشن وی
!!{lambda definable} ثابت ہو سکدے نیں۔ مثلاً، $f$ دی $g$
نال ترکیب $\lambd[x][F (G x)]$ نال !!{lambda defined} اے،
جے $f$ تے $g$ بالترتیب ٹرماں $F$ تے~$G$ نال
!!{lambda defined} ہون۔ پر جے فنکشن جزوی ہون، تاں ایس
وچ دِقّت اے۔ جدوں $g(x)$ غیر معین اے، اوہدا مطلب اے کہ
$G x$ دا کوئی نارمل روپ نہیں۔ زیادہ تر صورتاں وچ ایس دا
مطلب اے کہ $F (G
x)$ دا وی کوئی نارمل روپ نہیں، تے سانوں ایہو چاہیدا اے۔
پر جے $F$، $\lambd[x][\lambd[y][y]]$ ہووے، تاں $F (G x)$
دا نارمل روپ ($\lambd[y][y]$) موجود اے۔

ایہ دِقّت حل ہو سکدی اے۔ سارے جزوی بازگشتی فنکشوناں نوں
!!{lambda define} کرن دے ایہو جیہے طریقے موجود نیں جیہناں
وچ غیر معین قدراں نوں بغیر نارمل روپ والے ٹرم ظاہر کردے
نیں۔ پر ایہ طریقے اوہناں طریقیاں نالوں کجھ ودھ گُنجھلدار
نیں تے بدیہی طور اُتے گھٹ سمجھ آؤندے نیں جیہڑے اسیں عام
بازگشتی فنکشوناں لئی ورتے سن۔ ایتھے اسیں قضیہ ثبوت توں
بغیر لکھدے آں۔

\begin{thm}
  سارے جزوی بازگشتی فنکشن !!{lambda definable} نیں۔
\end{thm}

\end{document}
